\section{Conclusion}
\label{sec:conclusion}

We studied whether a conversion framework assembled for graph node classification transfers to graph-structured visual classification, where images ship as superpixel graphs rather than regular grids, and whether it can deliver an ANN-relative energy advantage. A linear global pooling operator preserves the exact conversion-error identity $\Bh-B=\RAEW$; a unified float64 forward pass verifies it to \IdentityMax. Within the transferred combination we found a genuine mechanistic redundancy: deterministic two-phase spiking already realises the target firing count, so the closed-form bias $\beta^\star=\theta/(2T)$ double-corrects the rate and must be removed. With that correction, membrane decoding, margin regularisation and a narrow low-timestep network, both 45\,nm arithmetic-energy conventions exceed $5\times$ across four mutually disjoint holdout groups and \TotalSeeds\ independent seeds, while all adapted single-point conversion baselines remain near chance. The accompanying cost is an explicit accuracy--energy frontier, including a negative scaling result that delimits the low-timestep regime. The broader lesson is that conversion discretises a computation along the neighbourhood as well as along time: a component that is correct for one operator can be wrong for another. Extending the analysis to dynamic vision graphs and replacing op-count estimates with hardware measurements follow directly from this boundary.
